\section{Stored Procedure dan Stored Function (Dasar)}

\textbf{Stored procedure} adalah kumpulan perintah SQL yang disimpan di database dan dapat dipanggil berulang kali. Stored procedure mengurangi pengiriman query dari aplikasi ke server, mengenkapsulasi logika bisnis di level database, dan meningkatkan keamanan karena user cukup diberi akses untuk memanggil procedure tanpa akses langsung ke tabel. \textbf{Stored function} mirip procedure tetapi mengembalikan satu nilai dan dapat digunakan di dalam query \cite{mysqlDocs}.

\begin{webcode}[language=SQL, caption={Stored procedure dan function}]
-- Stored Procedure: ambil produk berdasarkan kategori
DELIMITER //
CREATE PROCEDURE sp_produk_kategori(IN p_kategori VARCHAR(100))
BEGIN
  SELECT id, nama, harga FROM produk
  WHERE kategori_id = (
    SELECT id FROM kategori WHERE nama = p_kategori
  );
END //
DELIMITER ;

-- Memanggil procedure
CALL sp_produk_kategori('Elektronik');

-- Stored Function: hitung total produk
DELIMITER //
CREATE FUNCTION fn_total_produk() RETURNS INT
DETERMINISTIC
BEGIN
  DECLARE total INT;
  SELECT COUNT(*) INTO total FROM produk;
  RETURN total;
END //
DELIMITER ;

-- Menggunakan function
SELECT fn_total_produk() AS jumlah_produk;
\end{webcode}

